% Wersja źródłowa: audyt II, jednostka 12
\setcounter{section}{11}
\section{Klasyczne przestrzenie Banacha i ich przestrzenie sprzężone}
\label{sec:klasyczne-banach12}

\providecommand{\atlasref}[2]{#2}


\subsection{Przejście od ogólnych przestrzeni liniowo-topologicznych do norm}

Dotychczas prowadziliśmy rozważania dla przestrzeni liniowo-topologicznych,
przestrzeni lokalnie wypukłych i przestrzeni Frécheta. Teraz zawężamy
założenia, aby otrzymać konkretne wzory na normy, ciągłe funkcjonały
i przestrzenie sprzężone. Będziemy pracować nad $\K\in\{\R,\C\}$.
Niech $E'$ oznacza przestrzeń ciągłych funkcjonałów $\K$-liniowych na $E$.
Jeżeli $E$ jest unormowana, oznaczamy przez $\|\cdot\|_{E'}$ normę
operatorową na $E'$. Dla $1<p<\infty$ liczbę $q=p/(p-1)$ nazywamy
\emph{wykładnikiem sprzężonym} do $p$; dla $p=1$ przyjmujemy $q=\infty$.

\begin{definicja}[Przestrzeń unormowana i przestrzeń Banacha]
\label{def:banach12}
Przestrzenią unormowaną nazywamy parę $(E,\|\cdot\|)$, gdzie $E$
jest przestrzenią wektorową nad $\K$, a norma spełnia
$\|x\|=0\iff x=0$, $\|\alpha x\|=|\alpha|\|x\|$ oraz
$\|x+y\|\le\|x\|+\|y\|$. Przestrzeń unormowaną nazywamy
\emph{przestrzenią Banacha}, jeżeli jest zupełna względem metryki
$d(x,y)=\|x-y\|$.
\end{definicja}

\begin{stwierdzenie}[Norma operatora i przestrzeń sprzężona]
\label{prop:operator-norm12}
Niech $E$ będzie przestrzenią unormowaną i $F$ przestrzenią Banacha.
Wtedy przestrzeń $\mathcal L(E,F)$ wszystkich ciągłych operatorów
liniowych $E\to F$, z normą
$\|T\|=\sup_{\|x\|_E\le1}\|Tx\|_F$, jest przestrzenią Banacha.
W szczególności $E'=\mathcal L(E,\K)$ jest zawsze przestrzenią
Banacha, nawet gdy $E$ nie jest zupełna.
\end{stwierdzenie}
\begin{proof}
Ciągłość liniowego $T$ w zerze jest równoważna istnieniu $C\ge0$
takiego, że $\|Tx\|_F\le C\|x\|_E$: dla implikacji nietrywialnej
wystarczy przeskalować niezerowe $x$ do kuli, na której $T$ jest
ograniczony. Norma operatorowa jest zatem skończona i spełnia
aksjomaty normy.

Jeśli $(T_n)$ jest ciągiem Cauchy'ego w normie operatorowej,
to dla każdego $x\in E$ ciąg $(T_nx)$ jest Cauchy'ego w $F$, ponieważ
$\|T_nx-T_mx\|\le\|T_n-T_m\|\|x\|$.
Niech $Tx=\lim_nT_nx$. Granice zachowują liniowość. Ustalmy
$\varepsilon>0$ i wybierzmy $N$ tak, że
$\|T_n-T_m\|<\varepsilon$ dla $m,n\ge N$.
Po przejściu do granicy $m\to\infty$ otrzymujemy
$\|(T_n-T)x\|\le\varepsilon\|x\|$ dla $n\ge N$.
Stąd $T$ jest ciągły i $\|T_n-T\|\to0$.
\end{proof}

\subsection{Przestrzenie ciągowe}

\begin{definicja}[Przestrzenie $\ell^p$, $c_0$ i $\ell^\infty$]
\label{def:sequence-spaces12}
Dla $1\le p<\infty$ określamy
$\ell^p=\{x=(x_n)_{n\ge1}\in\K^{\N}:\sum_{n\ge1}|x_n|^p<\infty\}$
i $\|x\|_p=(\sum_{n\ge1}|x_n|^p)^{1/p}$.
Ponadto $\ell^\infty$ jest przestrzenią ciągów ograniczonych
z normą $\|x\|_\infty=\sup_{n\ge1}|x_n|$, a $c_0$ oznacza
podprzestrzeń ciągów zbieżnych do zera, z tą samą normą.
Przez $c_{00}$ oznaczamy przestrzeń ciągów o skończonym nośniku.
\end{definicja}

\begin{lemat}[Nierówność H\"oldera dla ciągów]
\label{lem:holder-seq12}
Dla $1<p<\infty$, $q=p/(p-1)$, $x\in\ell^p$ i $y\in\ell^q$
zachodzi
$\sum_{n\ge1}|x_ny_n|\le\|x\|_p\|y\|_q$.
Dla $p=1$, $q=\infty$, ta sama nierówność jest prawdziwa.
\end{lemat}
\begin{proof}
Dla $a,b\ge0$ nierówność Younga
$ab\le a^p/p+b^q/q$ wynika z wypukłości funkcji
$t\mapsto e^t$ (równoważnie, z minimalizacji
$t^p/p-tb$ względem $t\ge0$). Jeżeli obie normy są niezerowe,
stosujemy ją do $a=|x_n|/\|x\|_p$ i
$b=|y_n|/\|y\|_q$, sumujemy po $n\le N$,
a następnie przechodzimy do granicy $N\to\infty$.
Przypadki z normą zerową są oczywiste. Dla $p=1$
oszacowujemy bezpośrednio każdy wyraz przez
$|x_n|\|y\|_\infty$.
\end{proof}

\begin{twierdzenie}[Zupełność klasycznych przestrzeni ciągowych]
\label{thm:seq-banach12}
Dla każdego $1\le p\le\infty$ przestrzeń $\ell^p$ jest Banacha.
Przestrzeń $c_0$ jest domkniętą podprzestrzenią $\ell^\infty$,
a zatem także przestrzenią Banacha. Ponadto $c_{00}$ jest gęste
w $\ell^p$ dla $1\le p<\infty$ i w $c_0$, lecz nie jest gęste
w $\ell^\infty$.
\end{twierdzenie}
\begin{proof}
Dla $p=\infty$ ciąg Cauchy'ego $(x^{(k)})$ jest Cauchy'ego
jednostajnie względem współrzędnych. Istnieją granice
$x_n=\lim_kx_n^{(k)}$, a przejście do granicy w oszacowaniu
$\|x^{(k)}-x^{(m)}\|_\infty\le\varepsilon$ daje
$\|x^{(k)}-x\|_\infty\le\varepsilon$. Ciąg $x$ jest ograniczony.

Nierówność Minkowskiego dla skończonych sum ciągów
wynika z nierówności H\"oldera tak samo jak dla całek:
$\|u+v\|_p^p\le(\|u\|_p+\|v\|_p)\|u+v\|_p^{p-1}$
dla $p>1$, zaś dla $p=1$ jest natychmiastowa.

Dla $1\le p<\infty$ wybierzmy z ciągu Cauchy'ego
podciąg $(x^{(k_j)})$ spełniający
$\|x^{(k_{j+1})}-x^{(k_j)}\|_p\le2^{-j}$.
Niech $h_j=x^{(k_{j+1})}-x^{(k_j)}$.
Nierówność Minkowskiego dla sum skończonych daje
$\|\sum_{j=1}^N|h_j|\|_p\le\sum_{j=1}^N\|h_j\|_p\le1$.
Przez monotoniczność sum współrzędnych i ich granic
otrzymujemy $s=\sum_{j\ge1}|h_j|\in\ell^p$.
Szereg $\sum_jh_j$ jest więc bezwzględnie zbieżny
na każdej współrzędnej i określa $h\in\ell^p$;
ponadto norma jego ogona nie przekracza $\sum_{j\ge N}2^{-j}$.
Wobec tego $x^{(k_j)}\to x^{(k_1)}+h$ w $\ell^p$.
Cały ciąg Cauchy'ego ma tę samą granicę.

Jeśli $x^{(k)}\in c_0$ i $x^{(k)}\to x$ jednostajnie,
to dla $\varepsilon>0$ wybieramy $k$ z
$\|x-x^{(k)}\|_\infty<\varepsilon/2$, a następnie $N$
z $|x_n^{(k)}|<\varepsilon/2$ dla $n\ge N$.
Wtedy $x\in c_0$. Obcięcie ciągu do pierwszych $N$
współrzędnych daje zbieżność w $\ell^p$ dla $p<\infty$
i w $c_0$. W $\ell^\infty$ odległość stałego ciągu
$(1,1,\ldots)$ od $c_{00}$ wynosi co najmniej $1$.
\end{proof}

\begin{twierdzenie}[Przestrzenie sprzężone $\ell^p$ i $c_0$]
\label{thm:dual-sequences12}
Niech $1\le p<\infty$ i niech $q$ będzie wykładnikiem sprzężonym.
Każdy $F\in(\ell^p)'$ ma dokładnie jedną reprezentację
\[
 F(x)=\sum_{n=1}^\infty a_nx_n,\qquad x\in\ell^p,
 \quad a\in\ell^q,
\]
a odwzorowanie $a\mapsto F$ jest izometrycznym izomorfizmem
$\ell^q\cong(\ell^p)'$ (dla $p=1$ przyjmujemy $q=\infty$).
Ponadto każdy $F\in(c_0)'$ ma dokładnie jedną reprezentację
tego samego typu z $a\in\ell^1$ i $\|F\|=\|a\|_1$.
\end{twierdzenie}
\begin{proof}
Niech $e_n$ będzie $n$-tym wektorem standardowym i niech
$a_n=F(e_n)$. Dla $1<p<\infty$ i $N\ge1$ oznaczmy
$S_N=\sum_{n=1}^N|a_n|^q$. Jeśli $S_N>0$, wybieramy
$x_n=\overline{a_n}|a_n|^{q-2}S_N^{-1/p}$
dla $n\le N$ (a dla $a_n=0$ przyjmujemy $x_n=0$)
i $x_n=0$ dla $n>N$. Wtedy
$\|x\|_p=1$ i $F(x)=S_N^{1/q}$, ponieważ
$(q-1)p=q$. Zatem $S_N^{1/q}\le\|F\|$ dla każdego $N$.
Stąd $a\in\ell^q$ i $\|a\|_q\le\|F\|$.
Dla $p=1$ mamy bezpośrednio
$|a_n|\le\|F\|\|e_n\|_1=\|F\|$.

Ponieważ $c_{00}$ jest gęste w $\ell^p$ dla $p<\infty$,
na mocy ciągłości i Lematu~\ref{lem:holder-seq12}
$F(x)=\sum_na_nx_n$ dla każdego $x\in\ell^p$.
Nierówność H\"oldera daje $\|F\|\le\|a\|_q$,
a powyższe oszacowania dają równość.

Dla $F\in(c_0)'$ zdefiniujmy ponownie $a_n=F(e_n)$.
Dla każdego $N$ wektor $x\in c_{00}$ o współrzędnych
$x_n=\overline{a_n}/|a_n|$ przy $a_n\ne0$ i
$x_n=0$ przy $a_n=0$, $n\le N$, ma normę co najwyżej $1$.
Zatem $\sum_{n=1}^N|a_n|\le\|F\|$.
Wobec tego $a\in\ell^1$, a gęstość $c_{00}$ w $c_0$
oraz nierówność $|\sum_na_nx_n|\le\|a\|_1\|x\|_\infty$
kończą dowód wraz z identycznością norm.
Jednoznaczność w obu przypadkach wynika z podstawienia $x=e_n$.
\end{proof}

\begin{przyklad}[Dlaczego $(\ell^\infty)'$ jest większa niż $\ell^1$?]
\label{ex:linfty-dual-seq12}
Niech $u=(1,1,\ldots)\in\ell^\infty$. Na podprzestrzeni $c=c_0+\K u\subseteq\ell^\infty$
ciągów zbieżnych określmy $L(x)=\lim_{n\to\infty}x_n$.
Jest to ciągły funkcjonał liniowy o normie $1$.
Z twierdzenia Hahna--Banacha przedłużamy go do
$\widetilde L\in(\ell^\infty)'$ z tą samą normą.
Mamy $\widetilde L(e_n)=0$ dla każdego $n$, ale
$\widetilde L(u)=1$. Gdyby $\widetilde L$
był reprezentowany przez pewne $a\in\ell^1$ wzorem
$\widetilde L(x)=\sum_na_nx_n$, to z podstawienia $e_n$
wynikałoby $a_n=0$ dla wszystkich $n$, co przeczy
$\widetilde L(u)=1$.
\end{przyklad}

\subsection{Przestrzenie \texorpdfstring{$L^p$}{Lp} na przestrzeni z miarą skończoną}

\noindent\emph{Powiązania tematyczne w Bibliotece matematycznej:} \atlasref{ME-14::thm:riesz-fischer-14}{ME-14: twierdzenie Riesza--Fischera}; \atlasref{ME-15::thm:dual-lp-15}{ME-15: reprezentacja funkcjonałów na przestrzeniach funkcji całkowalnych}.


W tej podsekcji $(\Omega,\mathcal A,\mu)$ jest przestrzenią
z miarą dodatnią i skończoną, tj. $\mu(\Omega)<\infty$.
Rozważamy funkcje mierzalne o wartościach w $\K$, utożsamiane,
gdy są równe $\mu$-prawie wszędzie. Zapis
$\int_\Omega f(x)\,\mathrm d\mu(x)$ zawsze oznacza całkę
względem wskazanej miary. Jeżeli $\mu(\Omega)=0$, wszystkie
przestrzenie $L^p(\mu)$ są zerowe i twierdzenia pozostają prawdziwe.

\begin{definicja}[Przestrzenie $L^p(\mu)$]
\label{def:lp12}
Dla $1\le p<\infty$ niech $L^p(\mu)$ oznacza przestrzeń
klas funkcji $f$ spełniających
$\int_\Omega|f(x)|^p\,\mathrm d\mu(x)<\infty$, z normą
$\|f\|_p=(\int_\Omega|f(x)|^p\,\mathrm d\mu(x))^{1/p}$.
Przestrzeń $L^\infty(\mu)$ składa się z klas funkcji istotnie
ograniczonych z normą
$\|f\|_\infty=\operatorname*{ess\,sup}_{x\in\Omega}|f(x)|$.
\end{definicja}

\begin{twierdzenie}[H\"older, Minkowski i zupełność $L^p$]
\label{thm:lp-banach12}
Dla $1\le p\le\infty$ przestrzeń $L^p(\mu)$ jest Banacha.
Dla $1<p<\infty$ i $q=p/(p-1)$ zachodzi nierówność H\"oldera
$\int_\Omega|f(x)g(x)|\,\mathrm d\mu(x)\le\|f\|_p\|g\|_q$,
a dla $p=1$ analogiczny wzór z $q=\infty$.
Dla każdego $1\le p\le\infty$ zachodzi nierówność Minkowskiego
$\|f+g\|_p\le\|f\|_p+\|g\|_p$.
\end{twierdzenie}
\begin{proof}
Dowód H\"oldera powtarza argument z nierównością Younga
z Lematu~\ref{lem:holder-seq12}, po zastąpieniu sumy całką.
Dla $1<p<\infty$ najpierw $|f+g|^p\le2^{p-1}(|f|^p+|g|^p)$
zapewnia $f+g\in L^p$. Następnie nierówność Minkowskiego
otrzymujemy z $|f+g|^p\le(|f|+|g|)|f+g|^{p-1}$
i nierówności H\"oldera
i podzielenia przez $\|f+g\|_p^{p-1}$, o ile norma ta jest
niezerowa. Dla $p=1$ i $p=\infty$ jest to bezpośrednia
nierówność trójkąta.

Niech $(f_n)$ będzie ciągiem Cauchy'ego w $L^p$ dla
$1\le p<\infty$. Wybierzmy podciąg $(f_{n_k})$ taki, że
$\|f_{n_{k+1}}-f_{n_k}\|_p\le2^{-k}$.
Dla reprezentantów $h_k=f_{n_{k+1}}-f_{n_k}$,
Minkowski i twierdzenie o zbieżności monotonicznej dają
\[
 s(x)=\sum_{k\ge1}|h_k(x)|<\infty\quad\text{dla prawie każdego }x,
 \qquad \|s\|_p\le\sum_{k\ge1}2^{-k}.
\]
Szereg $\sum_kh_k$ zbiega więc prawie wszędzie do mierzalnej
funkcji $h\in L^p$, a normy ogonów są oszacowane przez
$\sum_{k\ge N}2^{-k}$. Stąd $f_{n_k}\to f_{n_1}+h$ w $L^p$.
Cauchy'ego całego ciągu zapewnia tę samą granicę.

Dla $p=\infty$ z warunku Cauchy'ego wybieramy podciąg
$(f_{n_k})$ z $\|f_{n_{k+1}}-f_{n_k}\|_\infty\le2^{-k}$.
Po usunięciu przeliczalnej sumy zbiorów zerowych
reprezentantów mamy punktową jednostajną zbieżność
szeregu różnic do istotnie ograniczonej funkcji.
Reszta dowodu przebiega jak dla $p<\infty$.
\end{proof}

\begin{lemat}[Gęstość funkcji prostych]
\label{lem:simple-dense12}
Dla $1\le p<\infty$ funkcje proste są gęste w $L^p(\mu)$.
Każda funkcja prosta jest kombinacją liniową funkcji charakterystycznych
zbiorów mierzalnych.
\end{lemat}
\begin{proof}
Dla rzeczywistej funkcji nieujemnej wybieramy rosnące funkcje
proste $s_n$ z $0\le s_n\uparrow f$ prawie wszędzie
(np. przez obcinanie wartości do $n$ i zaokrąglanie w dół
do wielokrotności $2^{-n}$). Wtedy
$|f-s_n|^p\le|f|^p$ i $|f-s_n|^p\to0$ prawie wszędzie.
Twierdzenie o zbieżności zmajoryzowanej daje
$\|f-s_n\|_p\to0$. Dla funkcji rzeczywistych ogólnych
i zespolonych rozdzielamy części dodatnie i ujemne
oraz części rzeczywiste i urojone.
\end{proof}

\begin{twierdzenie}[Przestrzeń sprzężona $L^p$ dla $1\le p<\infty$]
\label{thm:lp-dual12}
Niech $(\Omega,\mathcal A,\mu)$ będzie przestrzenią z miarą
skończoną, niech $1\le p<\infty$ i $q$ będzie wykładnikiem
sprzężonym. Dla każdego $F\in(L^p(\mu))'$ istnieje dokładnie
jedna klasa $g\in L^q(\mu)$ taka, że
\[
 F(f)=\int_\Omega f(x)g(x)\,\mathrm d\mu(x)
 \quad\text{dla wszystkich }f\in L^p(\mu).
\]
Ponadto $\|F\|=\|g\|_q$, więc
$(L^p(\mu))'\cong L^q(\mu)$ izometrycznie. W przypadku
zespolonym jest to identyfikacja liniowa przy przyjętej konwencji
wzoru bez sprzężenia $g$.
\end{twierdzenie}
\begin{proof}
Nierówność H\"oldera pokazuje, że dowolne $g\in L^q$
definiuje ciągły funkcjonał o normie co najwyżej $\|g\|_q$.
Dla danego $F\in(L^p)'$ definiujemy miarę zespoloną
(lub znakowaną w przypadku rzeczywistym)
$\nu(A)=F(\chi_A)$, $A\in\mathcal A$.
Funkcje charakterystyczne należą do $L^p$, ponieważ $\mu$ jest skończona.
Jeżeli $A_j$ są parami rozłączne, to
$\chi_{\bigcup_{j=1}^N A_j}\to\chi_{\bigcup_{j\ge1}A_j}$
w $L^p$ na mocy zbieżności monotonicznej; ciągłość $F$
daje przeliczalną addytywność $\nu$.
Dla dowolnego skończonego rozkładu $\Omega$ na zbiory
mierzalne $A_j$ wybierzmy liczby $\theta_j\in\K$ o module
nie większym niż $1$ tak, aby
$\theta_j\nu(A_j)=|\nu(A_j)|$. Wówczas
\[
 \sum_j|\nu(A_j)|
 =F\Bigl(\sum_j\theta_j\chi_{A_j}\Bigr)
 \le\|F\|\,\mu(\Omega)^{1/p}.
\]
Zatem $\nu$ ma skończoną wariację i jest absolutnie ciągła
względem $\mu$. Z twierdzenia Radona--Nikodyma dla miar
skończonych otrzymujemy $g\in L^1(\mu)$ takie, że
$\nu(A)=\int_Ag(x)\,\mathrm d\mu(x)$ dla $A\in\mathcal A$.
Przez liniowość i Lemat~\ref{lem:simple-dense12} trzeba jeszcze
wykazać, że $g\in L^q$.

Jeżeli $p=1$, to dla dowolnego $A\in\mathcal A$ stosujemy
nierówność normową do funkcji
$f=\chi_A\overline g/|g|$ na $\{g\ne0\}$,
a poza tym zbiorem przyjmujemy $f=0$. Z rozszerzenia wzoru
z funkcji prostych na funkcje ograniczone (przez jednostajną
aproksymację funkcjami prostymi) otrzymujemy
$\int_A|g|\,\mathrm d\mu\le\|F\|\mu(A)$.
Wybór $A=\{|g|>\|F\|+\varepsilon\}$ pokazuje,
że $\|g\|_\infty\le\|F\|$.

Jeżeli $1<p<\infty$, niech
$g_N=g\chi_{\{|g|\le N\}}$ i niech
$f_N=\overline{g_N}|g_N|^{q-2}$, z wartością $0$ w
punktach, gdzie $g_N=0$. Ponieważ $\mu$ jest skończona,
$f_N$ jest ograniczona, a więc należy do $L^p$.
Wzór reprezentacyjny obowiązuje dla niej, gdyż
jest ona granicą jednostajną funkcji prostych; wobec tego
\[
 \int_\Omega|g_N(x)|^q\,\mathrm d\mu(x)
 =F(f_N)\le\|F\|\left(\int_\Omega|g_N(x)|^q\,\mathrm d\mu(x)\right)^{1/p}.
\]
Jeśli całka po lewej stronie jest dodatnia, dzielimy przez
jej potęgę $1/p$ i otrzymujemy $\|g_N\|_q\le\|F\|$.
Zbieżność monotoniczna daje $g\in L^q$ i $\|g\|_q\le\|F\|$.
Dla $p=1$ wykazaliśmy analogiczną nierówność.
Teraz równość $F(f)=\int fg\,\mathrm d\mu$ wynika z gęstości
funkcji prostych i ciągłości obu stron na $L^p$.
Nierówność H\"oldera daje $\|F\|\le\|g\|_q$, czyli
równość norm. Jednoznaczność $g$ wynika z podstawiania
funkcji charakterystycznych zbiorów mierzalnych.
\end{proof}

\begin{uwaga}[Dlaczego ograniczenie $p<\infty$ jest istotne?]
\label{rem:linfty-dual12}
Przy $p=\infty$ przestrzeń sprzężona $(L^\infty(\mu))'$ jest
na ogół istotnie większa od $L^1(\mu)$. W przypadku miary
Lebesgue'a na $[0,1]$ wybierzmy parami rozłączne zbiory mierzalne
$A_n$ o dodatniej mierze, a $A=\bigcup_nA_n$. Operator
$P:L^\infty([0,1])\to\ell^\infty$ dany wzorem
$P(f)_n=\lambda(A_n)^{-1}\int_{A_n}f(x)\,\mathrm d\lambda(x)$
jest ciągły, liniowy i $\|P\|\le1$.
Jeżeli $\widetilde L$ jest przedłużeniem funkcjonału granicy
z Przykładu~\ref{ex:linfty-dual-seq12}, to $\Phi=\widetilde L\circ P$
spełnia $\Phi(\chi_{A_n})=0$ dla wszystkich $n$, lecz
$\Phi(\chi_A)=1$.
Gdyby $\Phi(f)=\int_{[0,1]}f(x)g(x)\,\mathrm d\lambda(x)$
dla pewnego $g\in L^1$, przeliczalna addytywność całki dałaby
$\Phi(\chi_A)=\sum_n\Phi(\chi_{A_n})=0$, sprzeczność.
\end{uwaga}

\subsection{Przestrzeń \texorpdfstring{$C(K)$}{C(K)} i twierdzenie Riesza--Markowa}

\noindent\emph{Powiązania tematyczne w Bibliotece matematycznej:} \atlasref{ME-21::thm:compact-riesz-21}{ME-21: twierdzenie Riesza--Markowa--Kakutaniego na przestrzeni zwartej}.


\begin{definicja}[Przestrzeń funkcji ciągłych i miary regularne]
\label{def:ck-measures12}
Niech $K$ będzie niepustą zwartą przestrzenią Hausdorffa.
Przez $C(K;\K)$ oznaczamy przestrzeń ciągłych funkcji
$K\to\K$ z normą $\|f\|_\infty=\max_{x\in K}|f(x)|$.
Skończoną miarę borelowską $\nu$ na $K$ nazywamy
\emph{regularną}, jeśli dla każdego borelowskiego $B\subseteq K$
i $\varepsilon>0$ istnieje zbiór domknięty $F$ i zbiór otwarty
$U$ takie, że $F\subseteq B\subseteq U$ oraz
$|\nu|(U\setminus F)<\varepsilon$, gdzie $|\nu|$ oznacza
miarę wariacji całkowitej. Zbiór takich miar ze znakiem
lub zespolonych oznaczamy przez $\mathcal M(K;\K)$.
\end{definicja}

\begin{stwierdzenie}[Zupełność $C(K)$]
\label{prop:ck-banach12}
Niech $\K\in\{\R,\C\}$. Dla każdej niepustej zwartej przestrzeni Hausdorffa $K$
przestrzeń $C(K;\K)$ jest przestrzenią Banacha.
\end{stwierdzenie}
\begin{proof}
Każda funkcja z $C(K;\K)$ jest ograniczona dzięki zwartości. Dla ciągu Cauchy'ego $(f_n)$ w normie supremum istnieje punktowa granica $f$, dzięki zupełności $\K$. Jeśli $|f_n(x)-f_m(x)|<\varepsilon/2$ dla $n,m\ge N$ i każdego $x$, to przechodząc do granicy przy $m\to\infty$ dostajemy $\|f_n-f\|_\infty\le\varepsilon/2$ dla $n\ge N$. Dla ustalonego $x$ wybieramy $n$ z $\|f_n-f\|_\infty<\varepsilon/3$, a następnie otoczenie $x$, na którym $|f_n(y)-f_n(x)|<\varepsilon/3$. Nierówność trójkąta daje $|f(y)-f(x)|<\varepsilon$. Stąd $f\in C(K;\K)$ i zbieżność w normie supremum.
\end{proof}

\begin{lemat}[Reprezentacja dodatniego funkcjonału]
\label{lem:riesz-positive12}
Niech $K$ będzie niepustą zwartą przestrzenią metryczną,
a $T:C(K;\R)\to\R$ dodatnim funkcjonałem liniowym,
tj. $T(f)\ge0$ dla każdej $f\ge0$.
Wówczas istnieje dokładnie jedna skończona regularna dodatnia miara
borelowska $\mu_T$ na $K$, dla której
$T(f)=\int_Kf(x)\,\mathrm d\mu_T(x)$ dla wszystkich $f\in C(K;\R)$.
\end{lemat}
\begin{proof}
Jest to dokładnie przypadek dodatniego funkcjonału na niepustej
zwartej przestrzeni metrycznej z twierdzenia
\atlasref{ME-21::thm:compact-riesz-21}{ME-21: twierdzenie Riesza--Markowa--Kakutaniego}.
W tym miejscu korzystamy z pełnego dowodu zawartego w teorii miary,
łącznie z konstrukcją miary, jej regularnością i jednoznacznością.
\end{proof}

\begin{twierdzenie}[Twierdzenie Riesza--Markowa na zwartej przestrzeni metrycznej]
\label{thm:riesz-markov12}
Niech $K$ będzie niepustą zwartą przestrzenią metryczną.
Dla każdego $F\in C(K;\K)'$ istnieje dokładnie jedna
regularna skończona miara borelowska $\nu\in\mathcal M(K;\K)$
taka, że
\[
 F(f)=\int_K f(x)\,\mathrm d\nu(x)
 \quad\text{dla wszystkich }f\in C(K;\K).
\]
Odwzorowanie $\nu\mapsto F$ jest izometrycznym izomorfizmem
$\mathcal M(K;\K)\cong C(K;\K)'$, przy normie
$\|\nu\|=|\nu|(K)$.
\end{twierdzenie}
\begin{proof}
Najpierw $\K=\R$. Dla $f\ge0$ definiujemy
$F^+(f)=\sup\{F(g):g\in C(K;\R),\ 0\le g\le f\}$.
Wartość supremum należy do $[0,\|F\|\|f\|_\infty]$.
Funkcja $F^+$ jest addytywna na stożku dodatnim.
Jedna nierówność wynika z wyboru $g_1\le f_1$ i
$g_2\le f_2$. Dla przeciwnej, jeśli
$0\le h\le f_1+f_2$, rozkładamy
$h=\min\{h,f_1\}+(h-\min\{h,f_1\})$;
oba składniki są ciągłe, nieujemne i ograniczone odpowiednio
przez $f_1$ i $f_2$. Po przejściu do supremum mamy
$F^+(f_1+f_2)=F^+(f_1)+F^+(f_2)$.
Rozszerzamy zatem $F^+$ na rzeczywisty funkcjonał liniowy
przez $F^+(f)=F^+(f_+)-F^+(f_-)$, gdzie
$f_+=\max\{f,0\}$, $f_-=\max\{-f,0\}$.
Jest on dodatni i ciągły; również $F^-=F^+-F$ jest dodatni,
ponieważ $F^+(f)\ge F(f)$ dla $f\ge0$.
Lemat~\ref{lem:riesz-positive12} daje dodatnie regularne miary
$\nu^+$ i $\nu^-$ reprezentujące odpowiednio $F^+$ i $F^-$.
Zatem $F(f)=\int_Kf\,\mathrm d(\nu^+-\nu^-)$.

Dla $\K=\C$ bierzemy części rzeczywistą i urojoną
ograniczenia $F$ do rzeczywistych funkcji ciągłych,
reprezentujemy je dwiema miarami rzeczywistymi i składamy
$\nu=\nu_1+i\nu_2$. Zespolona liniowość $F$ zapewnia,
że ten sam wzór działa na wszystkich $f\in C(K;\C)$.
Jednoznaczność wynika z jednoznaczności reprezentacji
części rzeczywistych i urojonych.

Dla dowolnej miary o skończonej wariacji mamy
$|\int f\,\mathrm d\nu|\le\|f\|_\infty|\nu|(K)$,
czyli $\|F\|\le|\nu|(K)$. W drugą stronę niech
$u=d\nu/d|\nu|$ będzie pochodną Radona--Nikodyma;
$|u|=1$ $|\nu|$-prawie wszędzie. Regularność miary
$|\nu|$ i własność aproksymacji funkcji charakterystycznych zbiorów
borelowskich funkcjami ciągłymi (przez zbiory domknięte
i otwarte oraz funkcje Urysohna) dają gęstość $C(K;\K)$
w $L^1(|\nu|)$. Aproksymując $\overline u$ funkcjami
ciągłymi o module nie większym niż $1$, otrzymujemy
$\|F\|\ge\lim|\int f_nu\,\mathrm d|\nu||=|\nu|(K)$.
To dowodzi izometrii.
\end{proof}

\begin{przyklad}[Miary Diraca i funkcjonały ewaluacji]
\label{ex:dirac12}
Dla $x_0\in K$ miara Diraca $\delta_{x_0}$ jest zdefiniowana
przez $\delta_{x_0}(B)=1$, gdy $x_0\in B$, i $0$ w przeciwnym
przypadku. Funkcjonał $F_{x_0}(f)=f(x_0)$ spełnia
$F_{x_0}(f)=\int_Kf\,\mathrm d\delta_{x_0}$ i ma normę $1$.
Jeśli $x_0\ne x_1$, funkcjonały ewaluacji są różne,
ponieważ metryka $K$ pozwala skonstruować funkcję ciągłą
rozdzielającą $x_0$ i $x_1$.
\end{przyklad}

\begin{uwaga}[Przygotowanie do przestrzeni dystrybucji]
\label{rem:bridge-distribution12}
Miara $\nu$ działa liniowo na funkcjach testowych przez całkowanie,
a miara Diraca działa przez ewaluację. Te same mechanizmy wystąpią
w przestrzeni dystrybucji $\mathcal D'(\Omega)$, jednak tam
przestrzeń funkcji testowych ma topologię granicy induktywnej,
a nie normę supremum. W szczególności nie wolno bez dodatkowych
argumentów utożsamiać wszystkich dystrybucji z miarami.
\end{uwaga}
